\appendix
\chapter{Questionnaire}\label{app:questionnaire}

Questionnaire used in the study. It includes four sections corresponding to math anxiety, self-efficacy, attitudes, and gender stereotypes (see \autoref{sec:Instruments} for details).

Both the English questionnaire and its Italian translation used in the study are reported below.

\newpage
\newgeometry{
  a4paper,
  inner=1.5cm,
  outer=1.5cm,
  top=1.5cm,
  bottom=1.5cm
}

\setlength{\headwidth}{\textwidth}
\fancyhfoffset{0pt}

%*****************************************
% ENGLISH VERSION
%*****************************************
\begin{center}
  { \textbf{English Version}}
\end{center}

\noindent
\begin{minipage}{0.65\textwidth}
  {\Large \textbf{QUESTIONNAIRE}}
\end{minipage}
\hfill
\begin{minipage}{0.30\textwidth}
  \raggedleft
  \textbf{Class:} \underline{\hspace{1cm}} \textbf{Age:} \underline{\hspace{0.8cm}}
\end{minipage}

\vspace{0.2cm}

% --- BOX CODE ---
\noindent
\fbox{
    \begin{minipage}{0.96\textwidth}
        \centering
        \vspace{0.1cm}
        \small
        \textbf{CODE}\\
        \footnotesize
        Do not write your name. Insert the letters/numbers indicated in the boxes below.\\
        \textit{Example: \textbf{Bi}anchi \textbf{Ma}rco, born in September (09) 2012 $\to$ BIMA0912}
        
        \vspace{0.1cm}
        \renewcommand{\arraystretch}{1.2}
        \setlength{\tabcolsep}{0pt} 
        
        \begin{tabular}{| *{8}{>{\centering\arraybackslash}p{1.4cm}|} }
            \hline
            \multicolumn{2}{|c|}{\footnotesize \textbf{SURNAME}} & 
            \multicolumn{2}{c|}{\footnotesize \textbf{NAME}} & 
            \multicolumn{2}{c|}{\footnotesize \textbf{MONTH}} & 
            \multicolumn{2}{c|}{\footnotesize \textbf{YEAR}} \\
            
            \multicolumn{2}{|c|}{\tiny (First 2 letters)} & 
            \multicolumn{2}{c|}{\tiny (First 2 letters)} & 
            \multicolumn{2}{c|}{\tiny (01-12)} & 
            \multicolumn{2}{c|}{\tiny (e.g. 12)} \\
            \hline
            \rule{0pt}{1cm} & & & & & & & \\
            \hline
        \end{tabular}
        \vspace{0.1cm}
    \end{minipage}
}

\vspace{0.1cm}
\noindent{\small \textbf{Gender:} $\bigcirc$ Female \quad $\bigcirc$ Male \quad $\bigcirc$ Other \quad $\bigcirc$ Prefer not to answer }
\vspace{0.1cm}
\hrule
% --- END HEADER ---

\vspace{0.1cm}

\noindent\textbf{Instructions:} Imagine being in the situations described below. Evaluate each situation in terms of how much anxiety you would feel, putting an \textbf{X} in the column that corresponds to your level of anxiety. There are no right or wrong answers.
\newline
\renewcommand{\arraystretch}{1.2} 
\setlength{\tabcolsep}{3pt} 

% TABLE 1
\noindent
\begin{tabularx}{\textwidth}{|X|c|c|c|c|c|}
\hline
\multicolumn{1}{|c|}{\textbf{Situation}} & \multicolumn{5}{c|}{\textbf{Level of anxiety}} \\
\cline{2-6}
\multicolumn{1}{|c|}{} & \textbf{\scriptsize Very little} & \textbf{\scriptsize Little} & \textbf{\scriptsize Moderate} & \textbf{\scriptsize Quite a lot} & \textbf{\scriptsize A lot} \\
\hline
Using multiplication tables and schemes to do math exercises &   &   &   &   &   \\
\hline
Thinking about the math test &   &   &   &   &   \\
\hline
Paying attention to the teacher solving a difficult operation on the board &   &   &   &   &   \\
\hline
Taking a written math test &   &   &   &   &   \\
\hline
Doing many difficult math exercises at home &   &   &   &   &   \\
\hline
Paying attention during math class &   &   &   &   &   \\
\hline
Watching a classmate solving a math exercise &   &   &   &   &   \\
\hline
Being questioned “by surprise” in math &   &   &   &   &   \\
\hline
Dealing with a new math topic &   &   &   &   &   \\
\hline
\end{tabularx}

\vspace{0.1cm}

\noindent\textbf{Instructions:} Evaluate how good you think you are at handling the different situations. Put an \textbf{X} in the column corresponding to your level. There are no right or wrong answers.
\newline

% TABLE 2
\noindent
\begin{tabularx}{\textwidth}{|X|c|c|c|c|c|}
  \hline
  \textbf{How good are you at...} & \textbf{\scriptsize Not at all} & \textbf{\scriptsize A little} & \textbf{\scriptsize Medium} & \textbf{\scriptsize Quite} & \textbf{\scriptsize Very} \\
  \hline
  Learning mathematics & & & & & \\
  \hline
  Solving math problems & & & & & \\
  \hline
  Mental calculation & & & & & \\
  \hline
  Solving operations in columns & & & & & \\
  \hline
  Multiplication tables & & & & & \\
  \hline
  Completing homework & & & & & \\
  \hline 
  Focusing on studying math without being distracted & & & & & \\
  \hline
  Committing to math study when you have other things to do & & & & & \\
  \hline
  Organizing school activities & & & & & \\
  \hline
  Being interested in school subjects & & & & & \\
  \hline
\end{tabularx}

\newpage

% TABLE 3
\noindent\textbf{Instructions:} Evaluate how much you agree with the statements below. Put an \textbf{X} in the box corresponding to your opinion. No right or wrong answers.
\newline
\noindent
\begin{tabularx}{\textwidth}{|X|c|c|c|c|c|}
  \hline
  \textbf{How much do you agree?} & \textbf{\scriptsize Not at all} & \textbf{\scriptsize A little} & \textbf{\scriptsize Medium} & \textbf{\scriptsize Quite} & \textbf{\scriptsize Very} \\
  \hline
  I like studying mathematics & & & & & \\
  \hline
  Mathematics is useful in everyday life & & & & & \\
  \hline
  I would rather not have math lessons & & & & & \\
  \hline
  I find mathematics interesting & & & & & \\
  \hline
  I think mathematics will be useful for my future & & & & & \\
  \hline
  Mathematics motivates me to do my best & & & & & \\
  \hline 
  I do not feel involved during math lessons & & & & & \\
  \hline
  I would like to go deeper into mathematics beyond school & & & & & \\
  \hline
\end{tabularx}
\vspace{1cm}
\newline
% TABLE 4
\noindent\textbf{Instructions:} Evaluate how much you agree with the statements below. Put an \textbf{X} in the box corresponding to your opinion. No right or wrong answers.
\newline

\noindent
\begin{tabularx}{\textwidth}{|X|c|c|c|c|c|}
  \hline
  \textbf{How much do you agree?} & \textbf{\scriptsize Not at all} & \textbf{\scriptsize A little} & \textbf{\scriptsize Medium} & \textbf{\scriptsize Quite} & \textbf{\scriptsize Very} \\
  \hline
  I worry that my classmates think I have less math ability because of my gender & & & & & \\
  \hline
  When I get a bad grade in math, I fear it confirms a stereotype about my gender & & & & & \\
  \hline
  I think girls should not deal with mathematics because it is not suitable for them & & & & & \\
  \hline
  Sometimes I think I should not be good at math because of my gender & & & & & \\
  \hline
  I worry that my math teacher may think I am not good because of my gender & & & & & \\
  \hline
  In mathematics, boys have a natural talent that girls cannot make up with effort & & & & & \\
  \hline 
  To achieve the same results, boys need to put in less effort & & & & & \\
  \hline
  Girls need to work harder than boys to achieve the same results & & & & & \\
  \hline
  Despite effort, girls remain on average less good at math than boys & & & & & \\
  \hline
\end{tabularx}

\vspace{2cm}
\centerline{\textit{Thank you for your participation!}}

%*****************************************
% ITALIAN VERSION
%*****************************************
\newpage
\begin{center}
  {\Large \textbf{Italian Version}}
\end{center}

\noindent
\begin{minipage}{0.65\textwidth}
    {\Large \textbf{QUESTIONARIO}}
\end{minipage}
\hfill
\begin{minipage}{0.30\textwidth}
    \raggedleft
    \textbf{Classe:} \underline{\hspace{1cm}} \textbf{Età:} \underline{\hspace{0.8cm}}
\end{minipage}

\vspace{0.2cm}

% --- BOX CODICE UNIFICATO ---
\noindent
\fbox{
    \begin{minipage}{0.96\textwidth}
        \centering
        \vspace{0.1cm}
        \small
        \textbf{CODICE}\\
        \footnotesize
        Non scrivere il tuo nome. Inserisci le lettere/numeri indicati nelle caselle qui sotto.\\
        \textit{Esempio: \textbf{Bi}anchi \textbf{Ma}rco, nato a Settembre (09) del 20\textbf{12} $\to$ BIMA0912}
        
        \vspace{0.1cm}
        \renewcommand{\arraystretch}{1.2}
        \setlength{\tabcolsep}{0pt} 
        
        % Tabella unica a 8 colonne
        \begin{tabular}{| *{8}{>{\centering\arraybackslash}p{1.4cm}|} }
            \hline
            \multicolumn{2}{|c|}{\footnotesize \textbf{COGNOME}} & 
            \multicolumn{2}{c|}{\footnotesize \textbf{NOME}} & 
            \multicolumn{2}{c|}{\footnotesize \textbf{MESE}} & 
            \multicolumn{2}{c|}{\footnotesize \textbf{ANNO}} \\
            
            \multicolumn{2}{|c|}{\tiny (Prime 2 lettere)} & 
            \multicolumn{2}{c|}{\tiny (Prime 2 lettere)} & 
            \multicolumn{2}{c|}{\tiny (01-12)} & 
            \multicolumn{2}{c|}{\tiny (es. 12)} \\
            \hline
            \rule{0pt}{1cm} & & & & & & & \\
            \hline
        \end{tabular}
        \vspace{0.1cm}
    \end{minipage}
}

\vspace{0.1cm}
\noindent{\small \textbf{Genere:} $\bigcirc$ Femmina \quad $\bigcirc$ Maschio \quad $\bigcirc$ Altro \quad $\bigcirc$ Preferisco non rispondere }
\vspace{0.1cm}
\hrule
% --- FINE INTESTAZIONE ---

\vspace{0.1cm}

\noindent\textbf{Istruzioni:} Immagina di trovarti nelle situazioni descritte qui sotto. Valuta ogni situazione in termini di quanta ansia proveresti, mettendo una \textbf{crocetta} nella colonna che corrisponde al tuo grado di ansia. Non ci sono risposte giuste o sbagliate.
\newline
\renewcommand{\arraystretch}{1.2} 
\setlength{\tabcolsep}{3pt} 

% TABELLA 1 - CORRETTA CON TABULARX (Niente resizebox)
\noindent
\begin{tabularx}{\textwidth}{|X|c|c|c|c|c|}
\hline
\multicolumn{1}{|c|}{\textbf{Situazione}} & \multicolumn{5}{c|}{\textbf{Grado di ansia}} \\
\cline{2-6}
\multicolumn{1}{|c|}{} & \textbf{\scriptsize Molto poca} & \textbf{\scriptsize Poca} & \textbf{\scriptsize Moderata} & \textbf{\scriptsize Abbastanza} & \textbf{\scriptsize Molta} \\
\hline
Usare le tabelline e gli schemi per svolgere esercizi di matematica &   &   &   &   &   \\
\hline
Pensare alla verifica scritta di matematica &   &   &   &   &   \\
\hline
Seguire con attenzione l'insegnante che risolve alla lavagna un'operazione difficile &   &   &   &   &   \\
\hline
Fare una verifica scritta di matematica &   &   &   &   &   \\
\hline
Svolgere per casa molti esercizi difficili di matematica &   &   &   &   &   \\
\hline
Seguire con attenzione la lezione di matematica &   &   &   &   &   \\
\hline
Seguire un compagno che risolve un esercizio di matematica &   &   &   &   &   \\
\hline
Essere interrogato “a sorpresa” in matematica &   &   &   &   &   \\
\hline
Affrontare un nuovo argomento di matematica &   &   &   &   &   \\
\hline
\end{tabularx}

\vspace{0.1cm}

\noindent\textbf{Istruzioni:} Valuta quanto pensi di essere bravo/a ad affrontare le diverse situazioni. Metti una \textbf{crocetta} nella colonna corrispondente. Non ci sono risposte giuste o sbagliate.
\newline

% TABELLA 2 - CORRETTA CON TABULARX
\noindent
\begin{tabularx}{\textwidth}{|X|c|c|c|c|c|}
  \hline
  \textbf{Quanto sei bravo/a...} & \textbf{\scriptsize Per niente} & \textbf{\scriptsize Poco} & \textbf{\scriptsize Medio} & \textbf{\scriptsize Abbastanza} & \textbf{\scriptsize Molto} \\
  \hline
  Nell'imparare la matematica & & & & & \\
  \hline
  Nel risolvere problemi matematici & & & & & \\
  \hline
  Nel calcolo a mente & & & & & \\
  \hline
  Nel risolvere operazioni in colonna & & & & & \\
  \hline
  Nelle tabelline & & & & & \\
  \hline
  Nel finire i compiti che ti sono stati assegnati per casa & & & & & \\
  \hline 
  Nel concentrarti nello studio della matematica senza farti distrarre & & & & & \\
  \hline
  Impegnarti nello studio della matematica quando hai altre cose interessanti da fare & & & & & \\
  \hline
  Organizzarti nelle attività scolastiche & & & & & \\
  \hline
  Interessarti nelle materie scolastiche & & & & & \\
  \hline
\end{tabularx}

\newpage


\noindent\textbf{Istruzioni:} Valuta quanto sei d'accordo con le affermazioni proposte di seguito mettendo una \textbf{crocetta} nella casella che corrisponde al tuo pensiero. Non ci sono risposte giuste o sbagliate.
\newline

% TABELLA 3 - CORRETTA CON TABULARX
\noindent
\begin{tabularx}{\textwidth}{|X|c|c|c|c|c|}
  \hline
  \textbf{Quanto sei d'accordo?} & \textbf{\scriptsize Per niente} & \textbf{\scriptsize Poco} & \textbf{\scriptsize Medio} & \textbf{\scriptsize Abbastanza} & \textbf{\scriptsize Molto} \\
  \hline
  Mi piace studiare matematica & & & & & \\
  \hline
  La matematica è utile nella vita quotidiana & & & & & \\
  \hline
  Preferirei non avere lezioni di matematica & & & & & \\
  \hline
  Trovo la matematica interessante & & & & & \\
  \hline
  Penso che la matematica sia utile anche per il mio futuro & & & & & \\
  \hline
  La matematica è una materia che mi motiva a dare il meglio & & & & & \\
  \hline 
  Non mi sento coinvolto/a durante le lezioni di matematica & & & & & \\
  \hline
  Vorrei approfondire la matematica oltre quello che facciamo a scuola & & & & & \\
  \hline
\end{tabularx}

\vspace{1cm}

\noindent\textbf{Istruzioni:} Valuta quanto sei d'accordo con le affermazioni proposte di seguito. Non ci sono risposte giuste o sbagliate.
\newline

% TABELLA 4 - CORRETTA CON TABULARX
\noindent
\begin{tabularx}{\textwidth}{|X|c|c|c|c|c|}
  \hline
  \textbf{Quanto sei d'accordo?} & \textbf{\scriptsize Per niente} & \textbf{\scriptsize Poco} & \textbf{\scriptsize Medio} & \textbf{\scriptsize Abbastanza} & \textbf{\scriptsize Molto} \\
  \hline
  Sono preoccupato/a che i miei compagni pensino che io abbia meno capacità matematiche a causa del mio genere & & & & & \\
  \hline
  Quando prendo un brutto voto in matematica temo che questo confermi lo stereotipo sul mio genere & & & & & \\
  \hline
  Penso che le ragazze non debbano occuparsi di matematica perchè non è adatta a loro & & & & & \\
  \hline
  Alcune volte penso che non dovrei essere bravo/a in matematica a causa del mio genere & & & & & \\
  \hline
  Mi preoccupo che la mia insegnante di matematica possa pensare che non sono bravo/a a causa del mio genere & & & & & \\
  \hline
  In matematica, i ragazzi hanno un talento naturale che le ragazze non possono compensare con l’impegno & & & & & \\
  \hline 
  Per ottenere gli stessi risultati in matematica i ragazzi devono impegnarsi di meno & & & & & \\
  \hline
  Le ragazze devono lavorare più dei ragazzi per ottenere gli stessi risultati in matematica & & & & & \\
  \hline
  Nonostante l’impegno, le ragazze restano mediamente meno brave in matematica dei ragazzi & & & & & \\
  \hline
\end{tabularx}

\vspace{2cm}
\centerline{\textit{Grazie per la partecipazione!}}

\restoregeometry
\setlength{\headwidth}{\textwidth}
\fancyhfoffset{0pt}
